\chapter{Ethical, Social, and Environmental Considerations}

\section{Dataset and Representation}

CIFAR-10 is a small historical web-image benchmark and is not representative of
the diversity, resolution, or context encountered in real deployments. The
system is unconditional and its $32\!\times\!32$ outputs are intended for
research comparison only. Generated images should be labelled as synthetic and
must not be presented as factual records or as representative of people,
communities, or real events.

The low resolution and absence of personal conditioning reduce identity and
privacy risks, but they do not remove the general risks associated with
generative imagery. Any extension to faces, private data, or deployment-facing
content would require a separate data-governance review, provenance controls,
acceptable-use guidance, and an appropriate model card.

\section{Responsible Reporting}

The report preserves unfavourable Mean Flow results, discloses unmatched
training budgets, and avoids claiming universal algorithm superiority. FID and
IS are reported with their sample count and limitations. The project therefore
treats negative and inconclusive evidence as part of the result rather than
selecting only favourable checkpoints or metrics.

The public repository contains source code and configuration files, while local
datasets, environments, generated pairs, results, and large checkpoints are
excluded. This separation reduces accidental publication of bulky artefacts or
machine-specific data and makes the boundary between reproducible source and
generated evidence explicit.

\section{Compute and Environmental Cost}

The canonical submitted training histories total approximately 8.26 GPU-hours,
excluding abandoned attempts, evaluation, and checkpoint generation. A carbon
estimate cannot be calculated reliably because power draw, utilisation, grid
intensity, and execution location were not logged. The report therefore states
the measured GPU time without converting it into an unsupported emissions
figure.

Future experimental logging should include GPU model, power draw, utilisation,
energy consumption, and location. The proposed adaptive-inference direction in
Chapter~9 is also environmentally relevant: allocating additional evaluations
only to difficult samples may reduce average computation when compared with a
fixed high-NFE policy.
